\chapter{Reproducibility and Evidence Boundary}
\label{app:reproducibility-evidence-boundary}

This appendix records the scope of the reproducibility evidence available for the thesis. The accompanying project archive contains checked result tables, a result manifest, checksums, a source packet, and a decision register. The compact reproducibility package indexes the causal-run, predictive-run, data, environment, source-version, archive-copy, and unresolved-gap records without duplicating the detailed manifest. Together, those records support the numerical traceability of the displayed results and preserve distinctions between admitted evidence, excluded artifacts, and unresolved limitations.

The available records do not establish a complete clean-checkout reproduction. The exact raw and processed data lineage, numbered producing configurations, predictive split and checkpoint-to-export lineage, producing source versions, and archive-copy history are incomplete. A repository requirements file or current wrapper can describe an implementation interface, but is not evidence that it was the exact environment or command that produced an archived result.

The two external ICU datasets also remain subject to their own access and governance conditions. The repository does not establish institutional approval numbers, an ethics-board determination, consent basis, or the final data-governance wording required for submission or any prospective use. These gaps do not negate the checked numerical record, but they limit reproducibility, administrative compliance, and any stronger clinical interpretation.

For these reasons, the thesis treats the available materials as an evidence and provenance boundary: they make the reported workflow more inspectable while leaving the recovery of complete per-run lineage, qualified clinical review, and authoritative institutional documentation as necessary future work.
